\documentclass[machines,article,submit,pdftex,moreauthors]{Definitions/mdpi}
\pdfinfoomitdate=1
\pdftrailerid{}
\pdfsuppressptexinfo=15
\graphicspath{{figures/}}
\setlength{\headheight}{20pt}
% Date-free manuscript copy: retain only journal identity and pagination.
\AtBeginDocument{%
  \renewcommand{\today}{}%
  \renewcommand{\cright}{}%
  \renewcommand{\leftcolDates}{}%
  \renewcommand{\leftcolUpdateCitation}{}%
  \fancyhf{}%
  \fancyhead[L]{\footnotesize\emph{\journalname}}%
  \fancyhead[R]{\footnotesize\thepage{} of \pageref*{LastPage}}%
  \fancypagestyle{plain}{%
    \fancyhf{}%
    \fancyhead[L]{\footnotesize\emph{\journalname}}%
  }%
}
\hreflink{}
\firstpage{1}
\makeatletter
\setcounter{page}{\@firstpage}
\makeatother
\pubvolume{1}
\issuenum{1}
\articlenumber{0}
\pubyear{}
\copyrightyear{}
\Title{Observer-Based Hybrid Position Holding for an Electrohydraulic Agricultural Hitch: A Two-Chamber Virtual-Prototype Stress Test}
\newcommand{\orcidauthorA}{0009-0005-1928-6218}
\Author{Chenyu Cui$^{1}$, Peng Wang$^{1}$ and Yucheng Zhang$^{1,}$*\orcidA{}}
\AuthorNames{Chenyu Cui, Peng Wang and Yucheng Zhang}
\address{%
$^{1}$ \quad Prospective Research Laboratory, Institute of Computing Technology, Chinese Academy of Sciences, Beijing 100190, China}
\corres{Correspondence: zhangyucheng@ict.ac.cn}
\datereceived{}
\daterevised{}
\dateaccepted{}
\datepublished{}
\abstract{Long-duration position holding in electrohydraulic agricultural hitches requires coordination between motion control and compensation for leakage and load changes. This study evaluates an observer-based hybrid controller in a nonlinear two-chamber virtual prototype with valve dynamics, internal leakage, friction, and a finite accumulator. The controller estimates velocity and implement load from noisy position and chamber-pressure measurements. A hysteretic supervisor with an 8 s minimum dwell selects fine trim, locking, or accumulator support. Across 20 frozen 1800 s runs, the controller satisfied the combined 10 mm holding and post-transition accuracy criteria in every run. Mean and worst hold drops were 1.85 and 2.10 mm, respectively; the worst post-transition error was 8.90 mm. A matched PID--accumulator baseline also satisfied the criteria, but required more mode switches and accumulator oil. Its lower tracking error reveals a tracking--holding trade-off. PID--lock and sliding-mode--lock exceeded the holding criterion after a load step. Tests with fivefold leakage, doubled sensor noise, and halved accumulator flow passed, whereas a +2 mm position bias produced a 10.61 mm worst error. The results support reproducible controller screening within the specified virtual model; they do not establish calibrated machine performance.}
\keyword{agricultural machinery; electrohydraulic actuator; two-chamber model; three-point hitch; position holding; interface observer; virtual prototype; hybrid supervision}
\begin{document}

\section{Introduction}
Rear hitches must accomplish two tasks with fundamentally different time scales: move an implement to a commanded position and maintain that position after motion has ended. Position holding is not merely the final portion of a tracking transient. It governs implement-depth retention, hydraulic energy consumption, and valve activity while the tractor operates under a quasi-static load. During this phase, cylinder and lock-valve leakage, load changes, temperature, sensor noise, and the available hydraulic storage can all produce a slow position drift. A controller optimized only for tracking can therefore generate unnecessary valve activity or fail to distinguish a persistent loss path from a transient disturbance.

The practical importance of this distinction is increasing as tractor hitches are integrated with automated depth control and coordinated traction functions. A recent review of tractor electrohydraulic hitch systems documents the breadth of position, draft, and mixed-control architectures used in agricultural machinery \cite{sun2022}. At the component level, the flow characteristics of a tractor electrohydraulic hitch valve have been modelled and experimentally validated \cite{kumar2021}; at the system level, coordinated hitch-depth and travel-speed control has been studied as an efficiency-oriented agricultural operation \cite{sun2025}. These studies motivate a clear representation of the hydraulic and supervisory mechanisms that govern a transition from commanded motion to holding.

Recent tractor research has also reported hitch-position estimation using an EKF--Transformer architecture \cite{kim2025} and robust draft/position control with differential and integral sliding modes \cite{liu2023}. More generally, electrohydraulic servo research provides established sliding-mode and adaptive-control designs \cite{bai2010,tang2011,jin2013}. These contributions establish useful control ingredients, but their performance evidence has different objectives from a long-hold stress test: many studies emphasize trajectory tracking, draft regulation, or a specific hardware configuration rather than the entry logic, retained leakage path, and storage-assisted response during an extended hold.

For a holding study, the evaluation protocol must be distinguished from any claim of physical performance. A small average drop can arise because a controller has not entered its hold mode, because the disturbance is unrealistically smooth, or because a noise-free internal pressure state is available to the controller. Conversely, a conservative velocity gate can prevent hold entry even when position tracking is stable. Reporting the entry criterion, time horizon, load disturbance, and parameter range is therefore as important as reporting position error. Simulation verification and physical validation address different questions \cite{sargent2013}: repeatable seeded noise can test a model's internal behavior, but cannot establish the performance of a hydraulic machine.

This work therefore investigates a narrow, simulation-based question: can an observer-based supervisor coordinate tracking, bounded valve trim, locking, and finite accumulator compensation while using only noisy position and chamber-pressure measurements? The goal is not to estimate the field performance of a particular hitch. Instead, the virtual prototype makes its loss paths, pressure interface, mode transitions, parameter classes, and stress conditions inspectable so that a reviewer can distinguish internal model consistency from physical validation.

\section{Related Work}
\subsection{Hitch control and electrohydraulic motion regulation}
Research on tractor hitches covers system development, coordination of agricultural operations, and actuator control. Reviews describe the sensing and hydraulic architectures used in position, draft, and mixed control \cite{sun2022}. Valve-flow studies provide component-specific experimental evidence \cite{kumar2021}, while coordinated depth and speed control links hitch commands to operating efficiency \cite{sun2025}. Tractor-specific sliding-mode control \cite{liu2023} and hitch-position estimation \cite{kim2025} address complementary regulation and sensing tasks. These sources establish the application context; their measurements are not used to calibrate the virtual prototype in this study.

Hydraulic servo models describe the coupling between flow, pressure, mechanical load, and displacement \cite{jelali2003}. Adaptive robust motion control of hydraulic actuators \cite{yao2000}, adaptive sliding-mode control with unknown parameters \cite{guan2008}, and nonlinear adaptive suspension control \cite{alleyne1995} provide established approaches to nonlinear regulation. Force control requires attention to the pressure--load relationship \cite{alleyne2000}, while backstepping designs can explicitly include valve dynamics \cite{choux2010}. Valve dead-zone compensation is another distinct source of motion-control improvement \cite{deng2017}. These methods inform the actuator-level context, but a comparison of tracking error alone does not evaluate long-duration hold entry, switching activity, and finite stored-fluid use.

Recent studies also illustrate the diversity of electrohydraulic control problems. Feng et al. combined disturbance observation with sliding-mode spool-position control \cite{feng2025}. Wang and Song studied clearance compensation using a Hopfield neural network \cite{wang2024}, and Hanaffi et al. evaluated optimized fractional-order PID position tracking through hardware-in-the-loop implementation \cite{hanaffi2026}. Their objectives and validation settings differ from the present 30 min holding protocol. The proposed supervisor is therefore evaluated against complete PID--lock, sliding-mode--lock, and PID--accumulator control chains rather than presented as a replacement for all actuator-level control methods.

\subsection{Observation, hydraulic storage, and mode supervision}
Disturbance-observer-based control \cite{chen2016} and active disturbance rejection \cite{han2009} provide a broader context for estimating unmeasured effects from controller-facing signals. Hydraulic leakage detection using an extended Kalman filter is a related diagnostic problem \cite{an2005}; it should be distinguished from reconstructing an internal coefficient from an exact virtual-model transition. Here, velocity and load estimates enter the V2 control law, whereas its leakage estimate is diagnostic only. The earlier reduced-order leakage proxy is retained separately to make this distinction explicit.

Low-speed behavior also depends on mechanical and fluid assumptions. Friction modeling and compensation are well-established topics \cite{wit1995,olsson1998}, and experimental work shows that an oil--air mixture's effective bulk modulus depends on its physical conditions \cite{gholizadeh2013}. Reviews of pump-controlled cylinder drives further show why actuator architecture matters when interpreting hydraulic efficiency \cite{ketelsen2019}. The present valve-controlled model uses researcher-defined friction, bulk-modulus, leakage, and accumulator inputs. Accordingly, reduced accumulator-oil use is reported as a controller-activity metric; it is not converted into an unmeasured system-energy saving.

Hybrid-system theory separates continuous dynamics from discrete mode transitions \cite{goebel2009}. Stability analyses of switched systems \cite{liberzon1999} and average dwell time \cite{hespanha1999} motivate explicit switching rules, but a numerical minimum dwell time alone does not prove closed-loop stability. In this study, hysteresis and an 8 s dwell constraint limit mode activity in the evaluated simulations. The boundedness argument applies to the legacy coefficient proxy; the V2 switched closed loop is assessed numerically rather than assigned a general stability guarantee.

\subsection{Scope of the present contribution}
The contribution is the coordination and evaluation of low-activity holding modes under a shared measurement interface. The primary plant maintains separate chamber-pressure states, finite accumulator storage, and explicit spool and pressure limits. All four control chains use the same observer and tracking bandwidth. Supervisor selection on five development seeds is separated from 20 frozen holdout runs, and adverse sensor-bias outcomes remain in the results. Following the distinction between verification and validation \cite{sargent2013}, the evidence concerns the specified virtual prototype. The deterministic stress grid tests selected parameter interactions; it does not estimate variance-based sensitivity indices or a physical failure probability \cite{saltelli2010}.

\section{Materials and Methods}
\subsection{Legacy reduced-order hitch model}
The earlier benchmark plant is a reduced-order pressure-difference virtual surrogate, not a two-chamber fluid-conservation model. Let $x$, $v$, and $d$ denote piston position, velocity, and virtual pressure difference. At each explicit 10 ms step, the implementation evaluates
\begin{equation}
 d_{k+1}=\Pi_{[0,p_{\mathrm{rel}}]}\left[d_k+\Delta t\left\{a_m(d_k^\star-d_k)-g_m\mathcal{L}(d_k,T_k;C_l,g_T)+\frac{\beta_e}{V}q_{\mathrm{acc},k}\right\}\right],
\end{equation}
where $d_k^\star=(F_{L,k}+kx_k+F_f)/A+G_p\operatorname{sat}_{[-1,1]}(u_k)$. The mode coefficient $a_m$ is $a_{\mathrm{track}}$ or $a_{\mathrm{hold}}$; $g_m$ is one in tracking and the declared retained-loss term in hold. $\mathcal{L}$ is the implementation-defined virtual loss operator, with its scalar $C_l$ multiplied linearly above 25 $^\circ$C by $g_T$. The accumulator term is zero outside accumulator-assisted hold. Thus, $G_p$, $a_m$, $g_m$, and $g_T$ are explicitly reduced-order coefficients, rather than valve or fluid data-sheet parameters. The mechanical state is updated by
\begin{equation}
 m\dot v=Ad-F_L-cv-F_f(v)-kx,\qquad \dot x=v.
\end{equation}
For controller interfaces only, paired virtual pressure outputs are reconstructed from $d$ and a bounded common pressure; they are not separately conserved chamber states. Nominal runs expose the noise-free virtual pressure difference, while the independent observability audit supplies controller-facing position, pressure, and load perturbations without changing the plant. This distinction makes the nominal pressure interface an explicit virtual-model assumption rather than a sensor-performance claim.

An accumulator is represented by a precharge pressure, gas volume, discharge valve, and recharge branch. Its discharge flow is
\begin{equation}
q_{acc}=u_{acc}K_{acc}\max(p_{acc}-d,0),\qquad
p_{acc}=\max\left(1.01d,\;p_{acc,0}\frac{V_{acc,0}}{V_{acc,0}+V_{out}}\right).
\end{equation}
This is a virtual compensation branch with a bounded recharge proxy. It is not a calibrated representation of a particular accumulator, lock valve, or hydraulic circuit.

\subsection{Parameter provenance and assumption register}
All 45 V2 plant, protocol, and supervisor parameters are exported directly from the frozen Python dataclasses to \texttt{paper1\_v2\_parameter\_manifest.csv}. Each row records its code identifier, numerical value, unit, source class, model role, and stress or selection coverage. The exporter is rerun before packaging so that the manifest cannot silently diverge from the executable defaults. Table~\ref{tab:provenance} summarizes the model-role classes. Every value is labelled as a researcher-defined theoretical virtual input; none is represented as a manufacturer data-sheet value or calibrated physical measurement.

\begin{table}[H]
\centering
\caption{Traceable V2 parameter roles exported from the frozen virtual prototype.}
\label{tab:provenance}
\resizebox{\textwidth}{!}{%
\begin{tabular}{lcl}
\toprule
Source class & Count & Interpretation \\
\midrule
Hydraulic flow/loss/pressure & 14 & chamber, valve, leakage, and pressure-limit inputs \\
Mechanical virtual plant & 7 & inertia, friction, stiffness, and stroke inputs \\
Stored-energy state/constraint & 8 & finite accumulator and related protocol inputs \\
Measurement interface & 4 & position/pressure noise and declared bias inputs \\
Controller, supervisor, and protocol & 12 & thresholds, dwell, horizon, load, and numerical settings \\
\bottomrule
\end{tabular}
}
\end{table}

\subsection{V2 two-chamber plant and interface observer}
The primary V2 plant replaces the reduced pressure-difference surrogate with two chamber-pressure states. Using the standard pressure-continuity structure for hydraulic servo systems \cite{jelali2003,yao2000}, the pressure equations for chamber volumes $V_A(x)$ and $V_B(x)$ are
\begin{equation}
 \dot p_A=\frac{\beta_e}{V_A(x)}(q_A+q_{acc}-A\dot x-q_l),\qquad
 \dot p_B=\frac{\beta_e}{V_B(x)}(A\dot x-q_B+q_l),
\end{equation}
where $q_A$ and $q_B$ are directional orifice flows from the commanded spool, $q_l=C_l(T)(p_A-p_B)$ is internal leakage, and $q_{acc}$ is limited by a finite gas--oil accumulator state. The spool follows a first-order lag with a dead zone. Mechanical motion uses $m\ddot x=A(p_A-p_B)-F_L-c\dot x-F_f(\dot x)-kx$. Relief pressure and stroke boundaries are enforced at every 10 ms step. All hydraulic coefficients are parameterized virtual inputs; none is a target-machine data-sheet or calibration value.

The controller receives noisy position and chamber-pressure measurements. An interface observer updates position rate and load using the measured innovation and pressure balance; a leakage estimate is retained only as a diagnostic. The true plant velocity, load, and leakage are never passed to the control law. The hybrid supervisor enters hold after a 1 s position/velocity gate and applies an 8 s minimum mode dwell. A pressure deficit above 0.30 MPa or a position error above 4 mm can request accumulator compensation, while a low-error, low-velocity condition requests lock. These settings were selected on five development seeds by requiring all paths to pass and then minimizing accumulator-oil use, mode switching, and hold drop in that order. They were fixed before the holdout runs.

\begin{figure}[H]
\centering
\includegraphics[width=0.98\linewidth]{paper1_v2_architecture.png}
\caption{Primary V2 architecture. All controllers use the same noisy position/chamber-pressure interface and observer; the proposed supervisor selects fine trim, lock, or finite-accumulator support.}
\label{fig:v2-architecture}
\end{figure}

\subsection{Assumption ranges and stress screens}
Table~\ref{tab:assumptions} lists selected values with dedicated stress screens. Except for position noise, the listed ranges are deterministic one-at-a-time settings rather than estimated probability distributions. They test whether the conclusions are locally fragile within the virtual model; they do not quantify uncertainty of a physical hitch.

\begin{table}[H]
\centering
\caption{Selected V2 virtual-prototype inputs and one-at-a-time stress settings. Full provenance and all nominal values are in \texttt{paper1\_v2\_parameter\_manifest.csv}; none is an identified physical measurement.}
\label{tab:assumptions}
\resizebox{\textwidth}{!}{%
\begin{tabular}{llll}
\toprule
Quantity & Nominal value & Stress settings & Role in model \\
\midrule
Moving mass $m$ & 1200 kg & $0.7m$, $1.3m$ & actuator inertia \\
Viscous coefficient $c$ & 2200 N s m$^{-1}$ & $0.7c$, $1.3c$ & mechanical damping \\
Hitch stiffness $k$ & 18 000 N m$^{-1}$ & $0.7k$, $1.3k$ & load-position relation \\
Cylinder leakage $C_l$ & $5.0\times10^{-14}$ m$^3$ s$^{-1}$ Pa$^{-1}$ & $5C_l$ & chamber-to-chamber loss \\
Position-noise SD & 1.5 mm & 0.5, 2.0$\times$ & measured position \\
Accumulator valve coefficient & $5.0\times10^{-8}$ m$^3$ s$^{-1}$ Pa$^{-1/2}$ & 0.5$\times$ at 14 kN & finite compensation branch \\
\bottomrule
\end{tabular}
}
\end{table}

\subsection{Legacy control architecture}
Figure~\ref{fig:architecture} shows the earlier reduced-order architecture retained as a historical benchmark. During the commanded motion, the controller combines a pressure-balance feedforward term with PID, DI-SMAC, or boundary-layer sliding-mode feedback. When the accumulator is inactive, the known command and load define a leakage-free pressure-rate prediction $\dot{\Delta p}^{(0)}_k$. The pressure-transition residual gives a non-negative coefficient candidate,
\begin{equation}
 C^*_{l,k}=\frac{A\max\{a_m(d_k^\star-d_k)-(d_{k+1}-d_k)/\Delta t,0\}}{g_m d_k[1+g_T\max(T_k-25,0)]},
\end{equation}
where $g_m=1$ in tracking and $g_m=0.08(1-u_{comp})$ in hold. The proxy applies an exponential update $\hat C_{l,k+1}=(1-\alpha)\hat C_{l,k}+\alpha C^*_{l,k}$ with $\alpha=0.08$ and does not update while the accumulator is discharging, because its virtual flow would confound a leakage-only residual. This is a model inversion diagnostic, not a physical leakage-identification method.

The hold supervisor is entered only after the position error is within 10 mm and the velocity gate is satisfied for 1 s. It applies hysteresis between fine trim, lock, and accumulator modes. The holding pressure target uses the measured or estimated implement load. Accumulator mode is selected by either the leakage-flow threshold or a measured pressure deficit above 0.08 MPa, so that a load step can invoke stored-fluid support without being misclassified as a leakage event. When the accumulator mode is active, its discharge valve is driven by pressure deficit and a bounded position-error integral; the main valve remains in bounded fine trim. For the observability audit, the pressure interface can apply a declared first-order measurement filter; its time constant is part of the evaluated controller interface, not a plant calibration. This separation avoids the opposing main-valve/accumulator actions observed in an earlier implementation.

\begin{figure}[H]
\centering
\includegraphics[width=0.96\linewidth]{paper1_architecture.png}
\caption{Dual-mode hitch-control architecture with a model-level pressure-transition leakage proxy.}
\label{fig:architecture}
\end{figure}

\subsection{Boundedness argument}
The analysis assumes bounded load and leakage-rate variation, bounded inversion error, a bounded valve command, and operation below the relief-pressure limit. For a bounded candidate error $\epsilon_k=C^*_{l,k}-C_l$, the observation error $\tilde C_l=C_l-\hat C_l$ obeys
\begin{equation}
|\tilde C_{l,k+1}|\le(1-\alpha)|\tilde C_{l,k}|+\alpha\bar\epsilon.
\end{equation}
Thus, for $0<\alpha<1$, the proxy is ultimately bounded by the virtual-model inversion error. This local result does not establish robustness to arbitrary sensor, valve, or accumulator errors, or stability of the complete switched system \cite{liberzon1999,hespanha1999}. The separate measurement audit quantifies selected controller-facing perturbations and retains the unfiltered-noise failure. For the sliding variable $s=\lambda(x_r-x)+v$, the boundary-layer controller gives a negative drift outside the boundary layer when the switching gain exceeds the disturbance and noise bounds. Finally, if residual hold leakage is bounded by $\bar q_l$, a conservative screening condition is
\begin{equation}
\Delta x_{1800}\le \frac{1800\bar q_l}{A}+\Delta x_{valve}+\Delta x_{structure}.
\label{eq:drop}
\end{equation}
Equation~\eqref{eq:drop} is a parameterized criterion, not proof of a physical-hardware requirement.

\subsection{Simulation protocol and statistics}
The legacy virtual prototype was simulated at 10 ms for 1800 s. The primary V2 protocol uses the same sample interval and horizon, but compares four complete control chains on the two-chamber plant: PID--lock, sliding-mode--lock, PID--accumulator, and observer hybrid. All four chains share the same controller-facing measurements, observer output, and tracking-stage bandwidth; their hold actions remain independent. Development seeds 4100--4104 were used only for the V2 dwell/hysteresis scan; frozen holdout seeds 4120--4139 were not used for selection. Position followed a smooth 0.15 m to 0.35 m reference over 8 s, the load changed from 9 to 14 kN at 60 s, and temperature rose from 25 to 50 $^\circ$C. The pass screen requires hold entry, maximum hold drop no greater than 10 mm, maximum post-transition absolute error no greater than 10 mm, and zero relief activation.

The outcome metrics were full-run RMSE, terminal position error, hold-entry time, 1800 s position drop, command energy, and hold-entry rate. The filtered velocity gate, raw velocity gate, and their conjunction were evaluated separately. Ablations removed load feedforward, accumulator closed-loop compensation, the leakage proxy, hysteresis, or the hold structure. A 14 kN step, leakage multipliers, and a 60 s cylinder-leakage step from one to eight times nominal were used as virtual stress tests.

\section{Results}
\subsection{V2 complete-chain comparison}
The V2 results are summarized in Table~\ref{tab:v2-lift}. The observer hybrid entered hold on all 20 frozen paths and passed the combined criterion on 20/20, with a mean hold drop of 1.85 mm and a worst hold drop of 2.10 mm. PID--lock and sliding-mode--lock both failed the 10 mm screen after the load step, with mean drops of 16.04 and 15.99 mm. The matched PID--accumulator chain passed 20/20 and achieved a lower transition RMSE, but its mean hold drop was 3.44 mm. Thus, the primary claim is a holding-efficiency trade-off against a credible active baseline, rather than a universal pass/fail separation.

To test whether the nominal V2 result depended on the original holdout seeds, an independent audit performed after the controller was frozen used 20 new nominal seeds (4200--4219) and separate ten-seed stress groups. The observer hybrid passed all 20 nominal runs, with a 1.89 mm mean hold drop, a 2.12 mm worst hold drop, and an 8.60 mm worst post-transition error. The groups with fivefold leakage, doubled noise, and halved accumulator flow each passed all ten runs; their worst hold drops were 2.26, 2.48, and 2.40 mm, respectively. These runs were not used for supervisor selection and extend the numerical evidence without changing the controller or the model's stated scope.

\begin{table}[H]
\centering
\caption{Frozen V2 two-chamber comparison on 20 independent 1800 s paths.}
\label{tab:v2-lift}
\scriptsize
\resizebox{\textwidth}{!}{%
\begin{tabular}{lrrrrr}
\toprule
Controller & Hold entry & Pass screen & Mean drop (mm) & Worst drop (mm) & Worst post-transition error (mm) \\
\midrule
PID--lock & 20/20 & 0/20 & 16.04 & 16.14 & 15.38 \\
Sliding-mode--lock & 20/20 & 0/20 & 15.99 & 16.11 & 16.62 \\
PID--accumulator & 20/20 & 20/20 & 3.44 & 3.87 & 7.43 \\
Observer hybrid & 20/20 & 20/20 & 1.85 & 2.10 & 8.90 \\
\bottomrule
\end{tabular}
}
\end{table}

With the fixed 8 s dwell rule, the observer hybrid produced a mean of 13.1 mode switches over 1800 s, compared with 26.8 for PID--accumulator. Mean accumulator-oil use was 1.02 mL for the observer hybrid and 25.54 mL for PID--accumulator. Their mean post-transition RMSE values were 1.22 and 0.25 mm, respectively. Reduced drift and control activity therefore came at the cost of higher tracking error. No run activated pressure relief.

\begin{figure}[H]
\centering
\includegraphics[width=0.98\linewidth]{paper1_v2_results.png}
\caption{Frozen 30 min V2 comparison. Bars show mean hold drop (left) and worst post-transition error (right); diamonds mark the worst hold-drop path. The dashed line is the declared 10 mm screen.}
\label{fig:v2-lift-results}
\end{figure}

\subsection{V2 interface and parameter stress boundaries}
In the original V2 stress groups, the observer hybrid passed all ten runs with fivefold internal leakage, doubled position and pressure noise, and halved accumulator flow. Their worst hold drops were 2.23, 2.67, and 1.98 mm, respectively. Accumulator support was essentially inactive in the nominal, leakage, and half-flow cases, but mean oil use increased to 61.98 mL under doubled sensor noise. This increase reveals the cost of noise-induced mode triggering. A $+2$ mm position bias did not prevent hold entry, but produced a worst post-transition error of 10.61 mm. The nominal accuracy result therefore depends on the stated measurement conditions. Unlike the legacy model's ideal-pressure interface, the V2 model makes this dependence explicit through noisy measurements and an observer.

\subsection{Legacy reduced-order benchmark}
For continuity with the earlier virtual-prototype study, this subsection retains the reduced-order benchmark and its associated sensitivity audits as a historical secondary analysis. It is not used to select or validate the primary V2 controller above; its purpose is to document how the earlier surrogate behaves and where its interface assumptions fail.

\subsubsection{Nominal tracking and holding}
Table~\ref{tab:main} gives the 20-seed results under the filtered-velocity definition. The controllers have similar RMSE because they share the pressure feedforward and hold supervisor. The approximately 200 mm dynamic maximum error is the initial 0.20 m commanded displacement and must not be interpreted as steady-state tracking error. The hybrid controller had 4.27 mm terminal error and 0.022 mm mean hold drop.

\begin{table}[H]
\centering
\caption{Results for the parameterized 1800 s test (mean $\pm$ standard deviation, $n=20$). Hold-entry rate uses the filtered velocity gate.}
\label{tab:main}
\resizebox{\textwidth}{!}{%
\begin{tabular}{lrrrrr}
\toprule
Controller & RMSE (mm) & Terminal error (mm) & Hold entry (s) & Hold drop (mm) & Entry rate \\
\midrule
PID & $9.179\pm0.004$ & $4.138\pm0.253$ & $8.207\pm0.116$ & $0.017\pm0.021$ & 100\% \\
DI-SMAC & $9.233\pm0.008$ & $4.288\pm0.277$ & $8.464\pm0.190$ & $0.019\pm0.023$ & 100\% \\
Boundary-layer SMC & $9.228\pm0.005$ & $4.254\pm0.266$ & $8.503\pm0.181$ & $0.014\pm0.020$ & 100\% \\
Dual-mode hybrid & $9.236\pm0.006$ & $4.267\pm0.296$ & $8.523\pm0.157$ & $0.022\pm0.027$ & 100\% \\
\bottomrule
\end{tabular}
}
\end{table}

The raw-velocity and dual-gate definitions produced zero hold entries in all 20 runs. This is a material result: filter delay is part of the present hold-entry behavior, so the filtered result cannot be treated as a direct safety statement for a real hitch.

\subsubsection{Leakage and load-step stress tests}
Figure~\ref{fig:leak} shows the lock-valve leakage-path sensitivity. The mean drop remained below 0.088 mm at ten times the assumed lock-path leakage. This margin only describes the low-leakage parameterization. It cannot replace temperature-dependent lock-valve measurements.

\begin{figure}[H]
\centering
\includegraphics[width=0.76\linewidth]{paper1_lock_leakage_sensitivity.png}
\caption{Sensitivity of the complete controller to the assumed lock-valve micro-leakage path ($n=20$).}
\label{fig:leak}
\end{figure}

The more discriminating ablation is the 14 kN load step (Figure~\ref{fig:ablation}). With load feedforward and accumulator closed-loop compensation, the mean drop was 1.75 mm. Removing the accumulator closed loop increased it to 20.85 mm, and removing load feedforward increased it to 96.59 mm. The results show that the claimed benefit is conditional on a correct load signal, the modeled pressure-deficit trigger, and available accumulator pressure.

\begin{figure}[H]
\centering
\includegraphics[width=0.76\linewidth]{paper1_load_step_ablation.png}
\caption{Load-step ablation for a 14 kN step. The dashed line is the assumed 10 mm, 30 min requirement ($n=20$).}
\label{fig:ablation}
\end{figure}

\subsubsection{Assumption-range and virtual leakage-step screens}
Figure~\ref{fig:assumption-proxy} summarizes the legacy model's one-at-a-time stress tests. The legacy nominal inputs are recorded in \texttt{paper1\_parameter\_manifest.csv}, separately from the V2 inputs in Table~\ref{tab:assumptions}. Across the tested mass, damping, stiffness, position-noise, and cylinder-leakage settings, every five-seed group entered hold and its mean 1800 s drop was below 0.060 mm. Under the 14 kN load step, accumulator-flow capacity had a larger effect: 0.5$\times$ and 1.5$\times$ the legacy nominal coefficient gave mean drops of 5.86 mm and 0.273 mm, respectively. These results describe model sensitivities, not confidence bounds for a physical hydraulic circuit.

For a virtual cylinder-leakage step at 60 s from one to eight times the nominal coefficient, the pressure-transition proxy reached 90\% of the new coefficient in $0.255\pm0.005$ s ($n=20$) and ended at $1.200\times10^{-11}$ m$^3$ s$^{-1}$ Pa$^{-1}$; a fixed nominal prior remained at $1.5\times10^{-12}$ m$^3$ s$^{-1}$ Pa$^{-1}$. Both cases entered hold in 20/20 runs and had the same 0.022 mm mean drop. Thus, the test demonstrates internal state-transition and supervisory-mode consistency, not an independent position-holding advantage or physical leakage identification.

\begin{figure}[H]
\centering
\includegraphics[width=0.94\linewidth]{paper1_assumption_and_proxy_stress.png}
\caption{Selected assumption-range screen (left, $n=5$) and virtual one-to-eight-times leakage step (right, $n=20$). Both panels are parameterized model tests, not calibrated-hitch evidence.}
\label{fig:assumption-proxy}
\end{figure}

\subsubsection{Temperature-load envelope and raw-velocity threshold}
The controller was next evaluated over a parameterized temperature-load envelope: constant oil temperatures of 25, 40, and 50 $^\circ$C were crossed with post-step loads of 7, 10, and 14 kN ($n=20$ per case). The filtered 3 mm s$^{-1}$ velocity gate entered hold in all nine cases. Figure~\ref{fig:envelope} shows that the highest mean 1800 s position drop was 1.93 mm (95\% CI: 1.74--2.11 mm) at 50 $^\circ$C and 14 kN. The result broadens the declared virtual operating envelope; its leakage and thermal inputs define internal model conditions rather than a physical temperature or load validation.

The same protocol was also evaluated with an unfiltered velocity signal. At raw thresholds of 3 and 5 mm s$^{-1}$, none of the 20 runs entered hold. Relaxing the raw threshold to 10 mm s$^{-1}$ produced 20/20 entries and a 0.031 mm mean drop (95\% CI: 0.016--0.047 mm). This sensitivity result identifies a candidate signal-processing design trade-off; it does not validate the original 3 mm s$^{-1}$ requirement or establish an acceptable hardware safety threshold.

\subsubsection{Transition-separated accuracy audit}
Because the commanded motion spans 200 mm, the all-time RMSE is dominated by the planned transition and is not a steady-state precision metric. We therefore recomputed accuracy after the 8 s reference transition. In the nominal virtual condition, the 20-seed mean post-transition RMSE was 4.25 mm, the mean 95th-percentile absolute error was 4.53 mm, and the mean per-seed maximum absolute error was 9.54 mm; the worst seed reached 10.61 mm. Doubling the position-noise standard deviation gave a mean maximum absolute error of 8.89 mm and a worst seed of 11.62 mm. These values support a statistical virtual accuracy description, but not a universal hard $\pm10$ mm guarantee.

\begin{figure}[H]
\centering
\includegraphics[width=0.94\linewidth]{paper1_operating_envelope.png}
\caption{Parameterized temperature-load envelope under the filtered velocity gate (left), and sensitivity of the unfiltered velocity gate to its tolerance (right), $n=20$ per condition.}
\label{fig:envelope}
\end{figure}

\subsubsection{Cartesian virtual stress surface}
To test interactions beyond the one-at-a-time tests, we evaluated the complete legacy controller on a predeclared Cartesian grid crossing mass (0.7 and 1.3 times nominal), cylinder leakage (1 and 5 times nominal), constant temperature (25 and 50 $^\circ$C), post-step load (7 and 14 kN), and accumulator-flow coefficient (0.5 and 1.0 times nominal). The 32 cases used the same 20 noise seeds as the legacy nominal experiment; none was used for tuning the V2 controller. All 640 runs entered hold and passed the 10 mm, 1800 s holding criterion. The worst case combined 0.7$\times$ mass, 1$\times$ leakage, 50 $^\circ$C, 14 kN, and 0.5$\times$ accumulator flow, with a mean drop of 7.99 mm (95\% CI: 7.77--8.21 mm). Figure~\ref{fig:combined-stress} describes coverage of specified virtual conditions, not a physical failure probability.

\begin{figure}[H]
\centering
\includegraphics[width=0.99\linewidth]{paper1_combined_stress_surface.png}
\caption{Predeclared Cartesian virtual stress surface for the complete hitch controller. Each cell summarizes 20 seeds. Panels fix mass ($m$) and leakage multiplier ($L$); rows combine temperature ($T$) and accumulator-flow multiplier ($Q$), and columns show the post-step load. The color scale is fixed at 0--10 mm, and every cell passes the 10 mm screen.}
\label{fig:combined-stress}
\end{figure}

\subsubsection{Measurement-interface and observer-mismatch audit}
The legacy nominal protocol supplies a noise-free virtual pressure signal. A separate 20-seed audit (seeds 2060--2079) therefore perturbed controller-facing measurements and observer coefficients while keeping the plant fixed. A $+2$ mm position bias did not prevent hold entry, but raised the worst steady-state error to 12.67 mm. Unfiltered 0.1 MPa Gaussian pressure noise caused instability after hold entry: in every run, the steady-state error reached the 500 mm virtual-stroke boundary. Table~\ref{tab:measurement-audit} reports this failure explicitly; hold entry alone is not counted as successful holding.

A declared 200 ms first-order pressure filter provided a restricted mitigation test. With 0.05 MPa pressure noise, all 20 runs entered hold and the worst steady error was 9.72 mm. The combined filtered case adds a $+2$ mm position bias, $+0.1$ MPa pressure bias, 0.05 MPa pressure noise, a $+0.5$ kN load-measurement bias, and conservative observer gain/rate mismatches; it entered hold in all 20 runs with an 8.97 mm worst steady error. These results define an interface-specific virtual boundary. They neither identify a hardware sensor budget nor convert the nominal result into a universal $\pm10$ mm guarantee.

\begin{table}[H]
\centering
\caption{Independent controller-facing measurement and observer audit ($n=20$ per row). All plant parameters remain fixed virtual inputs. The raw-noise row is an intentional retained failure boundary.}
\label{tab:measurement-audit}
\scriptsize
\setlength{\tabcolsep}{3pt}
\begin{tabularx}{\linewidth}{@{}X r r r r@{}}
\toprule
Case & Hold entry & Steady RMSE (mm) & Mean P95 abs. error (mm) & Worst steady error (mm) \\
\midrule
Nominal interface & 20/20 & 4.24 & 4.63 & 10.56 \\
$+2$ mm position bias & 20/20 & 4.25 & 4.63 & 12.67 \\
0.10 MPa raw pressure noise & 20/20 then unstable & 486.56 & 500.00 & 500.00 \\
0.05 MPa pressure noise, 200 ms filter & 20/20 & 4.69 & 5.55 & 9.72 \\
Combined filtered mismatch & 20/20 & 2.50 & 2.90 & 8.97 \\
\bottomrule
\end{tabularx}
\end{table}

\section{Discussion}
The primary V2 comparison does not support a claim that the observer hybrid dominates every metric. The matched PID--accumulator chain has lower post-transition RMSE and also passes the declared 10 mm screen. The observer hybrid instead reduces mean hold drop by 46.3\%, halves the mean mode-switch count, and reduces nominal accumulator-oil use by approximately 96\% relative to that active baseline. Its differentiated value is therefore the coordination of low-activity holding modes, not intrinsically superior trajectory tracking.

The central V2 modeling choices are the chamber volumes, valve coefficients, leakage paths, accumulator properties, friction, and pressure-measurement interface. The 20 holdout runs sample the specified measurement-noise process; they do not characterize a population of physical machines. Fivefold leakage and halved accumulator flow remained within the numerical criteria, while doubled noise increased accumulator use and mode switching. The +2 mm position-bias failure shows that a low simulated hold drop alone does not establish position accuracy when the measurement bias is unknown.

A separate legacy fine-trim scan tested whether the load-step result depended on the post-transition gains. Multiplying those gains by 1.5 reduced the mean 14 kN hold drop from 1.75 to 0.73 mm, but increased the nominal mean maximum post-transition error from 9.54 to 9.92 mm; the worst seed reached 11.60 mm. Multiplying the gains by two reduced the mean drop to 0.29 mm but produced a 13.04 mm worst nominal error. The legacy benchmark therefore retains its baseline gains and reports this trade-off. This scan did not select or retune the V2 controller.

The legacy measurement audit was not used to select gains or exclude unfavorable outcomes. It identifies a failure at the tested 0.1 MPa unfiltered pressure-noise level and a restricted mitigation with 0.05 MPa noise and a 200 ms filter. These model-specific results cannot define a hardware sensor specification. More generally, the study complements observer-based hydraulic control \cite{chen2016,feng2025} by exposing the holding supervisor's dependence on its measurement interface. It does not claim a new disturbance-observer theory or a complete switched-system stability proof \cite{goebel2009,liberzon1999}.

\section{Conclusions}
This study evaluated an observer-based hybrid holding controller in a two-chamber electrohydraulic virtual prototype. Across 20 frozen 30 min runs, it satisfied the specified 10 mm holding and post-transition accuracy criteria and limited the worst hold drop to 2.10 mm. Compared with a matched PID--accumulator baseline, it reduced mean hold drop, mode switching, and accumulator-oil use, but had a higher post-transition RMSE. PID--lock and sliding-mode--lock exceeded the holding criterion after the load step. Tests with fivefold leakage, doubled sensor noise, and halved accumulator flow passed, whereas a +2 mm position bias produced a 10.61 mm worst error. The contribution is a reproducible theoretical method for evaluating the trade-off between accuracy and control activity. All reported performance values apply to the specified virtual model; they do not establish calibrated machine performance or a universal $\pm10$ mm guarantee.

\authorcontributions{Conceptualization, C.C.; Methodology, C.C.; Software, C.C.; Formal analysis, C.C.; Visualization, C.C.; Writing---original draft preparation, C.C.; Resources, C.C.; Validation, P.W.; Investigation, C.C.; Writing---review and editing, Y.Z.; Supervision, Y.Z.; Project administration, Y.Z.; Funding acquisition, Y.Z.}
\funding{This work has received funding from the Science and Technology Project of the Ministry of Agriculture and Rural Affairs of China.}
\institutionalreview{Not applicable.}
\informedconsent{Not applicable.}
\dataavailability{The V2 virtual-prototype source code, reproduction scripts, parameter manifests, and frozen JSON outputs are provided with the submission as Supplementary Materials. The existing public project archive is available at \url{https://github.com/Fat-pig-Cui/agricultural-tractor-virtual-prototypes}; its version history is indexed by the Zenodo concept DOI \url{https://doi.org/10.5281/zenodo.22754509}. A new version containing the V2 supplementary files will be deposited before publication.}
\conflictsofinterest{The authors declare no conflicts of interest.}

\reftitle{References}
\begin{thebibliography}{30}
\bibitem{sun2022} Sun, X.; Lu, Z.; Song, Y.; Cheng, Z.; Jiang, C.; Qian, J.; Lu, Y. Development Status and Research Progress of a Tractor Electro-Hydraulic Hitch System. \emph{Agriculture} \textbf{2022}, \emph{12}, 1547. \url{https://doi.org/10.3390/agriculture12101547}.
\bibitem{kumar2021} Kumar, S.; Tewari, V.K.; Bharti, C.K.; Ranjan, A. Modeling, simulation and experimental validation of flow rate of electro-hydraulic hitch control valve of agricultural tractor. \emph{Flow Measurement and Instrumentation} \textbf{2021}, \emph{82}, 102070. \url{https://doi.org/10.1016/j.flowmeasinst.2021.102070}.
\bibitem{sun2025} Sun, X.; Song, Y.; Lu, Z.; Deng, X. Research on a Coordinated Control Method of Tractor Electro-Hydraulic Hitch Tillage Depth and Travel Speed Based on Optimal Overall Efficiency and Economic Performance. \emph{Agriculture} \textbf{2025}, \emph{15}, 2232. \url{https://doi.org/10.3390/agriculture15212232}.
\bibitem{kim2025} Kim, E.-K.; Han, T.-H.; Lee, J.-H.; Han, C.-W.; Lim, R.-G. Tracking accuracy evaluation of autonomous agricultural tractors via rear three-point hitch estimation using a hybrid model of EKF-Transformer. \emph{Agriculture} \textbf{2025}, \emph{15}, 1475. \url{https://doi.org/10.3390/agriculture15141475}.
\bibitem{liu2023} Liu, C.; Gu, J.; Du, X.; Liu, C.; Du, Y.; Mao, E. Differential and integral sliding mode adaptive control algorithm for draft and position integrated control of electro-hydraulic hitch in agricultural tractor. \emph{INMATEH Agricultural Engineering} \textbf{2023}, \emph{70}, 369--378. \url{https://doi.org/10.35633/inmateh-70-36}.
\bibitem{bai2010} Bai, Y.; Li, P. Adaptive fuzzy sliding mode control for electro-hydraulic position servo system. \emph{2010 Chinese Control and Decision Conference} 2010, 3249--3253. \url{https://doi.org/10.1109/CCDC.2010.5498613}.
\bibitem{tang2011} Tang, R.; Zhang, Q. Dynamic sliding mode control scheme for electro-hydraulic position servo system. \emph{Procedia Engineering} \textbf{2011}, \emph{24}, 28--32. \url{https://doi.org/10.1016/j.proeng.2011.11.2596}.
\bibitem{jin2013} Jin, B. Chattering inhibition of variable rate reaching law sliding mode control for electro-hydraulic position servo system. \emph{Journal of Mechanical Engineering} \textbf{2013}, \emph{49}, 163--169. \url{https://doi.org/10.3901/JME.2013.10.163}.
\bibitem{sargent2013} Sargent, R.G. Verification and validation of simulation models. \emph{Journal of Simulation} \textbf{2013}, \emph{7}, 12--24. \url{https://doi.org/10.1057/jos.2012.20}.
\bibitem{jelali2003} Jelali, M.; Kroll, A. \emph{Hydraulic Servo-systems: Modelling, Identification and Control}; Springer: London, UK, 2003. \url{https://doi.org/10.1007/978-1-4471-0099-7}.
\bibitem{yao2000} Yao, B.; Bu, F.; Reedy, J.; Chiu, G.T.-C. Adaptive robust motion control of single-rod hydraulic actuators: Theory and experiments. \emph{IEEE/ASME Transactions on Mechatronics} \textbf{2000}, \emph{5}, 79--91. \url{https://doi.org/10.1109/3516.828592}.
\bibitem{guan2008} Guan, C.; Pan, S. Adaptive sliding mode control of electro-hydraulic system with nonlinear unknown parameters. \emph{Control Engineering Practice} \textbf{2008}, \emph{16}, 1275--1284. \url{https://doi.org/10.1016/j.conengprac.2008.02.002}.
\bibitem{alleyne1995} Alleyne, A.; Hedrick, J.K. Nonlinear adaptive control of active suspensions. \emph{IEEE Transactions on Control Systems Technology} \textbf{1995}, \emph{3}, 94--101. \url{https://doi.org/10.1109/87.370714}.
\bibitem{alleyne2000} Alleyne, A.; Liu, R. A simplified approach to force control for electro-hydraulic systems. \emph{Control Engineering Practice} \textbf{2000}, \emph{8}, 1347--1356. \url{https://doi.org/10.1016/S0967-0661(00)00081-2}.
\bibitem{choux2010} Choux, M.; Hovland, G. Adaptive Backstepping Control of Nonlinear Hydraulic-Mechanical System Including Valve Dynamics. \emph{Modeling, Identification and Control} \textbf{2010}, \emph{31}, 35--44. \url{https://doi.org/10.4173/mic.2010.1.3}.
\bibitem{deng2017} Deng, W.; Yao, J.; Ma, D. Robust adaptive precision motion control of hydraulic actuators with valve dead-zone compensation. \emph{ISA Transactions} \textbf{2017}, \emph{70}, 269--278. \url{https://doi.org/10.1016/j.isatra.2017.07.022}.
\bibitem{feng2025} Feng, J.; Jia, Y.; Gao, W.; Cai, P.; Zhang, K.; Ai, C.; Wei, L. Sliding Mode Control of Spool Position for a Lead-Type Electro-Hydraulic Proportional Multi-Way Valve Based on Disturbance Observation. \emph{Machines} \textbf{2025}, \emph{13}, 774. \url{https://doi.org/10.3390/machines13090774}.
\bibitem{wang2024} Wang, T.; Song, J. Clearance Nonlinear Control Method of Electro-Hydraulic Servo System Based on Hopfield Neural Network. \emph{Machines} \textbf{2024}, \emph{12}, 329. \url{https://doi.org/10.3390/machines12050329}.
\bibitem{hanaffi2026} Hanaffi, M.S.; Othman, J.; Mohamaddan, S.; Mohd Faudzi, A.A.; Ghani, M.F.; Karim, M.H.F.A.; Ghazali, R. Hardware-in-the-Loop Implementation of Optimized Fractional-Order PID Control for Electro-Hydraulic Position Tracking. \emph{Machines} \textbf{2026}, \emph{14}, 1133. \url{https://doi.org/10.3390/machines14101133}.
\bibitem{chen2016} Chen, W.-H.; Yang, J.; Guo, L.; Li, S. Disturbance-Observer-Based Control and Related Methods---An Overview. \emph{IEEE Transactions on Industrial Electronics} \textbf{2016}, \emph{63}, 1083--1095. \url{https://doi.org/10.1109/TIE.2015.2478397}.
\bibitem{han2009} Han, J. From PID to Active Disturbance Rejection Control. \emph{IEEE Transactions on Industrial Electronics} \textbf{2009}, \emph{56}, 900--906. \url{https://doi.org/10.1109/TIE.2008.2011621}.
\bibitem{an2005} An, L.; Sepehri, N. Hydraulic Actuator Leakage Fault Detection using Extended Kalman Filter. \emph{International Journal of Fluid Power} \textbf{2005}, \emph{6}, 41--51. \url{https://doi.org/10.1080/14399776.2005.10781210}.
\bibitem{wit1995} Canudas de Wit, C.; Olsson, H.; \AA{}str\"om, K.J.; Lischinsky, P. A new model for control of systems with friction. \emph{IEEE Transactions on Automatic Control} \textbf{1995}, \emph{40}, 419--425. \url{https://doi.org/10.1109/9.376053}.
\bibitem{olsson1998} Olsson, H.; \AA{}str\"om, K.J.; Canudas de Wit, C.; G\"afvert, M.; Lischinsky, P. Friction Models and Friction Compensation. \emph{European Journal of Control} \textbf{1998}, \emph{4}, 176--195. \url{https://doi.org/10.1016/S0947-3580(98)70113-X}.
\bibitem{gholizadeh2013} Gholizadeh, H.; Bitner, D.; Burton, R.; Schoenau, G. Modelling and Experimental Validation of the Effective Bulk Modulus of a Mixture of Hydraulic Oil and Air. In \emph{ASME/BATH 2013 Symposium on Fluid Power and Motion Control}; 2013; V001T01A054. \url{https://doi.org/10.1115/FPMC2013-4493}.
\bibitem{ketelsen2019} Ketelsen, S.; Padovani, D.; Andersen, T.O.; Ebbesen, M.K.; Schmidt, L. Classification and Review of Pump-Controlled Differential Cylinder Drives. \emph{Energies} \textbf{2019}, \emph{12}, 1293. \url{https://doi.org/10.3390/en12071293}.
\bibitem{goebel2009} Goebel, R.; Sanfelice, R.G.; Teel, A.R. Hybrid dynamical systems. \emph{IEEE Control Systems Magazine} \textbf{2009}, \emph{29}, 28--93. \url{https://doi.org/10.1109/MCS.2008.931718}.
\bibitem{liberzon1999} Liberzon, D.; Morse, A.S. Basic problems in stability and design of switched systems. \emph{IEEE Control Systems Magazine} \textbf{1999}, \emph{19}, 59--70. \url{https://doi.org/10.1109/37.793443}.
\bibitem{hespanha1999} Hespanha, J.P.; Morse, A.S. Stability of switched systems with average dwell-time. In \emph{Proceedings of the 38th IEEE Conference on Decision and Control}; 1999; Volume 3, pp. 2655--2660. \url{https://doi.org/10.1109/CDC.1999.831330}.
\bibitem{saltelli2010} Saltelli, A.; Annoni, P.; Azzini, I.; Campolongo, F.; Ratto, M.; Tarantola, S. Variance based sensitivity analysis of model output. Design and estimator for the total sensitivity index. \emph{Computer Physics Communications} \textbf{2010}, \emph{181}, 259--270. \url{https://doi.org/10.1016/j.cpc.2009.09.018}.
\end{thebibliography}
\end{document}
